\section{Mechanization and Refinement Chain}
\label{sec:mechanization}

Each guarded operation is assembled from three proof layers and a linking step.

\paragraph{Numerical theorem.}
Rocq functional models describe the sequence of rounded operations.  Flocq
theorems relate the represented output to the exact real operation and, where
applicable, establish the output representation relation.  FLX--FLT bridge
lemmas or direct bounded-format arguments discharge the exponent-range issues
described in Section~\ref{sec:ideas}.

\paragraph{Condition bridge.}
The operation-specific condition packages finiteness, representation, ordering,
and range hypotheses.  The numerical development proves a corollary from that
condition to the result consumed by the VST wrapper proof.

\paragraph{C refinement and guarded composition.}
A VST refinement proof establishes that the C kernel implements the binary64
functional model.  A separate body proof establishes the check specification
\[
  b=\mathsf{true}\Longrightarrow\operatorname{cond}(x).
\]
The wrapper proof follows its two branches.  On success, ordinary Hoare
sequencing supplies the check postcondition as the condition needed at the
internal call; after the kernel call, the numerical theorem supplies the
postcondition.  On failure, the specification of the error path makes an unsafe
kernel call unreachable.

The check and operation translation units are packaged as Verified Software
Units.  Linking verifies that imported and exported function specifications
match across translation-unit boundaries.  CompCert's correctness theorem then
applies to the linked Clight program along the selected compilation path
\cite{Appel2014}.

The trusted and assumed boundary must consequently be read in layers.  The
machine-checked results rely on the Rocq kernel and the stated assumptions of
the imported proof libraries.  The executable interpretation additionally
relies on CompCert and the platform behavior assumed by the C and IEEE~754
interfaces.  The runtime checks themselves are program components covered by
VST proofs; builds produced through a different compiler are a distinct
deployment choice and should not be described as the same verified binary.
